\documentclass[shortabstract,english,mgr]{thesis}

\usepackage[utf8]{inputenc}

%%%%% DANE DO STRONY TYTUŁOWEJ
\polishtitle  {Klasyfikacja pewnej klasy grafów kwantowych \fmlinebreak nad
zespolonymi macierzami $ 3 \times 3 $}
\englishtitle {Classification of a class of quantum graphs \fmlinebreak over
complex $ 3 \times 3 $ matrices}
\polishabstract {Choć nieskierowane i bezpętlowe grafy kwantowe na
    $M_2(\mathbb{C})$ doczekały się w ostatnich latach pełnej
    klasyfikacji, przejście do wyższych wymiarów stanowi zasadniczy
    przeskok strukturalny. W niniejszej pracy podejmujemy naturalny kolejny
    krok: klasyfikację tego typu grafów na $M_3(\mathbb{C})$. Dokonujemy
    pełnej klasyfikacji bezpętlowych, nieskierowanych grafów kwantowych
    generowanych przez uogólnione macierze Pauliego na $M_3(\mathbb{C})$,
    wykazując, że tworzą one pięć klas z podziałem ze względu na liczbę
    krawędzi, są wierzchołkowo przechodnie, a zarazem kwantowo
    izomorficzne z grafami klasycznymi. Ostatecznie dowodzimy, że
    przejrzysta klasyfikacja dostępna w wymiarze drugim załamuje się w
    przypadku $M_3(\mathbb{C})$ z powodu skomplikowanej struktury orbit
    pod wpływem działania grupy $PU(3)$, co ujawnia nieskończone rodziny
nieizomorficznych struktur.}
\englishabstract{While undirected and
    loop-free quantum graphs on $ M_2(\mathbb{C}) $ have been fully
    classified in recent years, moving to higher dimensions presents a
    profound sructural leap. In this thesis we investigate the natural
    next step: the classification of such graphs on $ M_3(\mathbb{C})
    $. We completely classify the loop-free, undirected quantum graphs
    generated by the generalized Pauli matrices on $ M_3(\mathbb{C}) $,
    demonstrating that they can be divided into five classes based on
    their number of edges, are vertex-transitive, and quantum
    isomorphic to classical derived graphs. Finally, we establish that
    the clean classification available in dimension two breaks down in $
    M_3(\mathbb{C}) $ due to a complex orbit structure under the $
    PU(3) $ action, revealing infinite families of non-isomorphic
structures.} \author         {Filip Rabiega}        % Autor \advisor
%       {Dr. Mariusz Tobolski} % Promotor
\date           {August 29th, 2026}  % Data złożenia pracy
\transcriptnum  {330273}               % Numer indeksu
%%%%%

%%%%% WLASNE DODATKOWE PAKIETY
\usepackage{graphicx,listings,amsmath,amssymb,amsthm,amsfonts,tikz,tikz-cd,float}
\usetikzlibrary{arrows.meta, positionand, calc, babel}

%%%%% CYTOWANIA
\usepackage[backend=biber]{biblatex}
\addbibresource{cits.bib}

%%%%% WŁASNE DEFINICJE I POLECENIA
\theoremstyle{definition} \newtheorem{definition}{Definition}[section]
\theoremstyle{plain} \newtheorem{theorem}{Theorem}[section]
\theoremstyle{plain} \newtheorem{proposition}{Proposition}[section]
\theoremstyle{plain} \newtheorem{lemma}{Lemma}[section]
\theoremstyle{plain} \newtheorem{fact}{Fact}[section]
\theoremstyle{remark} \newtheorem{corollary}{Corollary}[section]
\theoremstyle{remark} \newtheorem{remark}{Remark}[section]
\theoremstyle{remark} \newtheorem{example}{Example}[section]
\renewcommand \qedsymbol {\ensuremath{\square}}
%%%%%

\begin{document}

%%%%% ZASADNICZA CZĘŚĆ PRACY

% General Tikz Settings
\tikzset{
    base node/.style={circle, draw, minimum size=0.85cm,
    font=\sffamily\bfseries\normalsize, thick},
    derived node/.style={circle, draw, minimum size=0.72cm, inner
    sep=0.75pt, font=\sffamily\small\bfseries, thick, fill=white},
    edge label/.style={font=\sffamily\small, fill=white, inner sep=1.25pt},
}

% Colors used consistently for voltage 0 / 1 / 2 throughout this section
\colorlet{volt0}{blue}
\colorlet{volt1}{red}
\colorlet{volt2}{green!55!black}

% Draw the 3 base nodes on a small equilateral triangle
\newcommand{\drawbase}{
    \node[base node] (u) at (0,0) {\( u \)};
    \node[base node] (v) at (3.4,0) {\( v \)};
    \node[base node] (w) at (1.7,2.94) {\( w \)};
}

% Step 1: invisible coordinates for the 3x3 grid, spaced out to
% reduce edge crowding.
% Edges are drawn against these bare coordinates; the visible node
% circles are drawn
% afterwards (see \drawfibernodes) so that edges never paint over a
% node's fill or label.
\newcommand{\drawfibers}{
    \foreach \name/\x in {u/0, v/3.8, w/7.6} {
        \coordinate (\name0) at (\x, 5.1);
        \coordinate (\name1) at (\x, 2.55);
        \coordinate (\name2) at (\x, 0);
    }
}

% Step 2: the visible fiber labels and node circles, drawn last so
% they sit on top of
% every edge.
\newcommand{\drawfibernodes}{
    \foreach \name/\x in {u/0, v/3.8, w/7.6} {
        \node[font=\sffamily\bfseries\small, text=gray] at (\x, 6.2)
        {Fiber \name};
        \node[derived node, draw=volt0, text=volt0] at (\name0) {\(\name,0\)};
        \node[derived node, draw=volt1, text=volt1] at (\name1) {\(\name,1\)};
        \node[derived node, draw=volt2, text=volt2] at (\name2) {\(\name,2\)};
    }
}

% Draw the undirected intra-fiber loops (lifts of base loops form a
% C_3 in each fiber)
\newcommand{\drawderivedloops}{
    \foreach \name in {u, v, w} {
        \draw[-, thick, black, opacity=0.85] (\name0) to (\name1);
        \draw[-, thick, black, opacity=0.85] (\name1) to (\name2);
        \draw[-, thick, black, opacity=0.85, bend right=40]
        (\name0) to (\name2);
    }
}

\newcommand{\drawlift}[3]{
    % #1 = source fiber, #2 = target fiber, #3 = voltage
    \ifnum#3=0
    \draw[->, >=stealth, thick, volt0, opacity=0.7,
    shorten >=2.2pt, shorten <=2.2pt]
    (#10) to (#20);
    \draw[->, >=stealth, thick, volt0, opacity=0.7,
    shorten >=2.2pt, shorten <=2.2pt]
    (#11) to (#21);
    \draw[->, >=stealth, thick, volt0, opacity=0.7,
    shorten >=2.2pt, shorten <=2.2pt]
    (#12) to (#22);
    \fi
    \ifnum#3=1
    \draw[->, >=stealth, thick, volt1, opacity=0.6,
    shorten >=2.2pt, shorten <=2.2pt]
    (#10) to (#21);
    \draw[->, >=stealth, thick, volt1, opacity=0.6,
    shorten >=2.2pt, shorten <=2.2pt]
    (#11) to (#22);
    \draw[->, >=stealth, thick, volt1, opacity=0.6,
    shorten >=2.2pt, shorten <=2.2pt]
    (#12) to (#20);
    \fi
    \ifnum#3=2
    \draw[->, >=stealth, thick, volt2, opacity=0.6,
    shorten >=2.2pt, shorten <=2.2pt]
    (#10) to (#22);
    \draw[->, >=stealth, thick, volt2, opacity=0.6,
    shorten >=2.2pt, shorten <=2.2pt]
    (#11) to (#20);
    \draw[->, >=stealth, thick, volt2, opacity=0.6,
    shorten >=2.2pt, shorten <=2.2pt]
    (#12) to (#21);
    \fi
}

\newcommand{\drawliftgrey}[3]{
    % #1 = source fiber, #2 = target fiber, #3 = voltage
    \ifnum#3=0
    \draw[->, >=stealth, thick, darkgrey, opacity=0.7,
    shorten >=2.2pt, shorten <=2.2pt]
    (#10) to (#20);
    \draw[->, >=stealth, thick, darkgrey, opacity=0.7,
    shorten >=2.2pt, shorten <=2.2pt]
    (#11) to (#21);
    \draw[->, >=stealth, thick, darkgrey, opacity=0.7,
    shorten >=2.2pt, shorten <=2.2pt]
    (#12) to (#22);
    \fi
    \ifnum#3=1
    \draw[->, >=stealth, thick, darkgrey, opacity=0.6,
    shorten >=2.2pt, shorten <=2.2pt]
    (#10) to (#21);
    \draw[->, >=stealth, thick, darkgrey, opacity=0.6,
    shorten >=2.2pt, shorten <=2.2pt]
    (#11) to (#22);
    \draw[->, >=stealth, thick, darkgrey, opacity=0.6,
    shorten >=2.2pt, shorten <=2.2pt]
    (#12) to (#20);
    \fi
    \ifnum#3=2
    \draw[->, >=stealth, thick, darkgrey, opacity=0.6,
    shorten >=2.2pt, shorten <=2.2pt]
    (#10) to (#22);
    \draw[->, >=stealth, thick, darkgrey, opacity=0.6,
    shorten >=2.2pt, shorten <=2.2pt]
    (#11) to (#20);
    \draw[->, >=stealth, thick, darkgrey, opacity=0.6,
    shorten >=2.2pt, shorten <=2.2pt]
    (#12) to (#21);
    \fi
}

% One-line color legend, reused above every derived graph
\newcommand{\voltagelegend}{%
    \begin{tikzpicture}[baseline]
        \node[font=\sffamily\small] at (0,0)
        {\textcolor{volt0}{\textbf{---}} voltage 0 \quad
            \textcolor{volt1}{\textbf{---}} voltage 1 \quad
        \textcolor{volt2}{\textbf{---}} voltage 2};
    \end{tikzpicture}%
}

\chapter{Introduction}

\section{Motivation}
Quantum graphs is a topic that has seen a resurgence in interest in
recent years.

\section{Scope}

\section{Acknowledgements}
I would like to thank my supervisor Mariusz Tobolski for his unwavering support
and patience.

\chapter{Preliminaries}

\section{Group theory}

In the following section we will introduce some group theoretic notions that
be used throughout this thesis. We assume that the Reader is familiar with
the notion of a group.

% This is from Dummit & Foote's Algebra, page 167
\begin{definition}[Dual of a group]
    Let \( G \) be a group. The \emph{dual group} \( \hat{G} \) is the
    set of all homomorphisms from \( G \) into the multiplicative group
    of roots of unity in \( \mathbb{C} \) with the group operation defined
    as
    \[
        \left( \varphi \psi \right)(g) = \varphi(g) \psi(g)
    \]
    for all \( \varphi, \psi \in \hat{G} \) and \( g \in G\), where the latter
    multiplication takes place in \( \mathbb{C} \).
\end{definition}

\begin{remark}
    We will often call the elements of the dual group
    \emph{characters}. If a group is
    isomorphic to its complement it will be called \emph{self-dual}.
\end{remark}

The following proposition will be used throughout the thesis, as almost
all of the groups we will be dealing with will be finite and abelian.

\begin{proposition}
    If \( G \) is finite and abelian, then it is self-dual.
\end{proposition}

\begin{proof}
    First we will prove the statement for finite cyclic groups. Let \( G =
    \left\langle g \right\rangle \), \( \operatorname{ord}(g) = n \). We will
    prove that the dual is cyclic of order \( n \). Let \( \chi\colon
        G \rightarrow
    \mathbb{C} \) be a character. Since the group is cyclic, the character is
    entirely determined by the value it takes on the generator \( g \). Since
    \( g^{n} = e \), obviously \( \chi(g)^{n} = 1 \), so \( \chi(g)^{n} \in
        \mu_{n} = \left\{ \exp\left( 2k \pi i / n \right) \colon k =
            0, 1, \dots, n-1
    \right\}\). For each possible choice of \( \chi(g) \in \mu_{n} \) we get a
    different character. It is then easy to see that \( \chi\colon G \rightarrow
    \mathbb{C} \) such that \( \chi(g) = \exp\left( 2 \pi i / n \right) \)
    generates this group, i.e. \( \hat{G} = \left\{ \chi, \chi^{2},
    \dots, \chi^{n-1}, \chi^{n} = e_{\hat{G}} \right\} \).

    Now since every finite abelian group can be written as a finite product
    of finite cycles, all we need to do is check whether the statement
    holds for a product of two groups \( A, B \) that are themselves
    self-dual, i.e. \( \widehat{A \times B} \cong \widehat{A} \times
    \widehat{B} \). This is self-evident as every character
    \( \chi \in \widehat{A \times B} \) is determined by its restrictions
    to \( A \times \{e_{B}\} \) and \( \left\{ e_{A} \right\} \times B \), since
    \[
        \chi(a, b) = \chi((a, e_{B})(e_{A}, b)) = \chi(a, e_{B})\chi(e_{A}, b).
    \]
    Thus \( \chi \leftrightarrow (\chi|_{A}, \chi|_{B})
    \) is an isomorphism.
\end{proof}

\section{\(C^{*}\)-algebras}

In the following section we will introduce some notions stemming from the
theory of \( C^{*} \)-algebras. We assume the Reader is familiar with
the notion of a \(C^{*}\)-algebra. For reference, see \textcite{blackadar}.

Multiplication in a \(C^{*}\)-algebra \( A \) is a bilinear map
\[
    m \colon A \times A \to A.
\]
By the universal property of the tensor product, \( m \) extends uniquely
to a linear map
\[
    m^{*} \colon A \otimes A \to A.
\]
Throughout this thesis, we will use the same notation \( (m) \) for both
the bilinear multiplication and its unique linear extension, relying
on the context to distinguish between the two.

\begin{definition}[Linear functional]
    A \emph{linear functional} on a \(C^{*}\)-algebra \( A \) is a
    linear function \( \psi\colon A \rightarrow \mathbb{C} \).
\end{definition}

\begin{remark}
    We will call linear functionals simply as \emph{functionals}.
\end{remark}

The following notions will be useful for defining an inner product on
\(C^{*}\)-algebras. Positiveness and faithfulness are all we need to
define a norm on a finite-dimensional \(C^{*}\)-algebra equipped with
a suitable functional.

\begin{definition}[Positive functional]
    A functional \( \psi \) on a \(C^{*}\)-algebra \( A \) is called
    \emph{positive}, if
    \[
        \psi(a^{*}a) \geq 0
    \]
    for any \( a \in A \).
\end{definition}

\begin{definition}[Faithful functional]
    A positive functional \( \psi \) on a \(C^{*}\)-algebra \( A \) is
    called \emph{faithful}, if
    \[
        \psi(a^{*}a) = 0 \Leftrightarrow a = 0
    \]
    for any \( a \in A \).
\end{definition}

% TODO: What literature?
\begin{remark}
    In the literature, the functional \( \psi \) is often called a \emph{state}.
    % For example see \textcite{gromada22}.
\end{remark}

The following notion is akin to the notion of a group action on a set.

\begin{definition}[Group action on a \(C^{*}\)-algebra]
    By \emph{group action} of a group \( G \) on a \(C^{*}\)-algebra \(
    A \) we mean a group homomorphism
    \[
        \alpha\colon G \rightarrow \operatorname{Aut}(A).
    \]
    We will write \( g \cdot a \) instead of \( \alpha(g)(a) \) then the action
    in question is clear.
\end{definition}

\subsection{Crossed-product \(C^{*}\)-algebra}

% TODO: rewrite this subsection

In this section we will introduce the notion of a crossed-product
\(C^{*}\)-algebra. We will do so in a limited manner, only to the
extent to which the theory of such products is necessary for our
purposes. For a more general reference we refer the Reader to
\textcite{blackadar}.

The following definitions are based off of \textcite{sierakowski09}.

\begin{definition}[\( C^{*} \)-dynamical system]
    A \( C^{*} \)-dynamical system is a triple \( (A, G, \alpha) \),
    where \( A \) is a \(C^{*}\)-algebra, \( G \) is a group, and \(
    \alpha\colon G \rightarrow \operatorname{Aut}(A) \) is a group
    action on \( A \).
    For a \( C^{*} \) dynamical system \( (A, G, \alpha) \) with \( G
    \) finite we define
    \[
        C_{c}(G, A, \alpha) = \left\{ f\colon G \rightarrow A \colon
        \operatorname{supp}(f) \text{ is finite}\right\},
    \]
    i.e. the space of all finitely supported functions on \( G \) with values
    in \( A \).
\end{definition}

\( C_{c}(G, A, \alpha) \) is equipped with both a product and a \( *
\)-operation in such a way that the action becomes "inner", i.e. \( g
\cdot a = u_{g} a u^{*}_{g} \) for all \( g \in G \), \( a \in A \).
More precisely we define those two operations in the following way:
\[
    ab = \sum_{g,h \in G}
\]

\begin{definition}[Covariant representation]
    Fix a dynamical system \( (A, G, \alpha) \). A \emph{covariant
    representation} \( (\pi, u, H) \) of \( (A, G, \alpha) \) consist of a
    unitary representation \( u\colon G \rightarrow B(\mathcal{H}) \)
    of \( G \) and
    a representation \( \pi\colon A \rightarrow B(\mathcal{H}) \) of
    \( A \) such
    that \( u(g) \pi(a) u(g)^{*} = \pi(g \cdot a) \) for every \( g \in G \)
    and \( a \in A \).
\end{definition}

\begin{definition}[Full \(C^{*}\)-algebra norm on \( C_{c}(G, A, \alpha) \)]
    The full \(C^{*}\)-algebra norm on \( C_{c}(G, A, \alpha) \) is given by
    \[
        \left\| \cdot \right\| = \sup \left\| (\pi \times u)(\cdot) \right\|,
    \]
    where the supremum is taken over all covariant representations \(
    (\pi, u, H) \)
    of \( (A, G, \alpha) \).
\end{definition}

\begin{definition}[Crossed product \(C^{*}\)-algebra]
    The crossed product of a \(C^{*}\)-algebra and a discrete group \( G \) is
    the completion of \( C_{c}(G, A, \alpha) \) with respect to the full norm.
    We denote the crossed product by \( \tilde{B}
    \rtimes_{\hat{\alpha}} \hat{G} \).
\end{definition}

\chapter{Quantum sets and quantum graphs}

In this chapter we will introduce the notions of quantum sets and quantum
graphs. They were first considered by...

\section{Quantum sets}

Let \( B \) be a \(C^*\)-algebra equipped with a faithful positive functional
\( \psi \). By the Gelfand-Naimark theorem, \( B \) can be realized concretely
as a closed \(*\)-subalgebra of the bounded linear operators \( B(\mathcal{H})
\) on some Hilbert space \( \mathcal{H} \). When \( B \) is finite-dimensional,
this abstract framework simplifies considerably: we can identify \( B \) with
its own Hilbert space completion \( L^2(B, \psi) \) by defining the inner
product directly via the functional: \[ \langle x, y \rangle = \psi(x^* y)
\quad \forall x, y \in B. \]

The following definition introduces one of the central notions used
in this thesis. It is taken directly from \textcite{tobolski26}.

\begin{definition}[Quantum set, \textcite[Definition~2.1]{tobolski26}]
    A quantum set \( (B, \psi) \) is a pair consisting of a finite-dimensional
    \(C^{*}\)-algebra \( B \) and a faithful positive functional \( \psi\colon B
    \rightarrow \mathbb{C} \) such that
    \[
        mm^{*} = \operatorname{Id}_{B}
    \]
    holds for the multiplication map \( m\colon B \otimes B \rightarrow B
    \) and its adjoint \( m^{*} \) given by
    \[
        \left\langle m(a \otimes b), c \right\rangle = \left\langle a
        \otimes b, m^{*}(c) \right\rangle_{\psi \otimes \psi}
    \]
    for all \( a, b, c \in B \).
\end{definition}

\begin{remark}
    We will call any functional with the property given in the
    definition above a \emph{quantum set functional} on \( B \).
\end{remark}

\begin{remark}
    Recall that in general, a \emph{comultiplication} on a
    \(C^{*}\)-algebra is a \(*\)-homomorphism
    \[
        \Delta\colonA\to A\otimes A
    \]
    such that the following diagram commutes:
    \[
        \begin{tikzcd}
            A \arrow[r,"\Delta"] \arrow[d,"\Delta"']
            & A\otimes A \arrow[d,"{\Delta\otimes\mathrm{id}}"] \\
            A\otimes A \arrow[r,"{\mathrm{id}\otimes\Delta}"']
            & A\otimes A\otimes A
        \end{tikzcd}
    \]
    Note that in general a comultiplication is \emph{not} unique. In our
    setting it is uniquely determined by the additional structure and the
    Riesz representation theorem for Hilbert spaces.
\end{remark}

\begin{remark}
    Since in our setup the comultiplication is given by the scalar product
    and regular multiplication \( m \), we will denote it by \( m^{*} \).
\end{remark}

% TODO: Prove that the examples are correct (do some calculations)
\begin{example}
    Every finite set \( X \) gives rise to a quantum set \( (C(X), \psi_{X}) \)
    with \( \psi_{X}\colon C(X) \rightarrow \mathbb{C} \), \(
        \psi_{X}(f) = \sum_{x
    \in X} f(x) \) for any \( f \in C(X) \).
    Let \( X \) be a finite set and \( C(X) \) be the \(*\)-algebra of
    complex-valued functions on \( X \), where the involution is given by
    pointwise complex conjugation: \( f^*(x) = \overline{f(x)} \).

    To show that \( \psi_X \) is positive, we must verify that
    \( \psi_X(f^* f) \geq 0 \) for all \( f \in C(X) \). For any function
    \( f \), the product \( f^* f \) evaluates pointwise to:
    \[
        (f^* f)(x) = \overline{f(x)} f(x) = |f(x)|^2.
    \]
    Applying the functional yields:
    \[
        \psi_X(f^* f) = \sum_{x \in X} |f(x)|^2.
    \]
    Thus the functional is positive.

    To prove faithfulness, suppose \( \psi_X(f^* f) = 0 \) for some \( f
    \in C(X) \). This implies:
    \[
        \sum_{x \in X} |f(x)|^2 = 0.
    \]
    Because each term \( |f(x)|^2 \) is strictly non-negative, a sum of
    such terms can only equal zero if every individual term is zero.
    Thus, we must have:
    \[
        |f(x)|^2 = 0 \quad \forall x \in X,
    \]
    which directly implies \( f(x) = 0 \) for all \( x \in X \). Therefore,
    \( f \) is the zero function (\( f = 0 \)), establishing that \(
    \psi_X \) is faithful.

    Let \( \{e_x\}_{x \in X} \) be the standard orthonormal basis of
    characteristic functions on \( X \). Recall the definitions of the
    multiplication map \( m \) and its adjoint \( m^* \):
    \begin{align*}
        m(f \otimes g) &= fg, \\
        m^*(e_x) &= e_x \otimes e_x.
    \end{align*}

    We wish to show that the composition \( mm^*\colon C(X)
    \rightarrow C(X) \) is the identity operator (\( \text{Id} \)). Let us
    evaluate the action of \( m m^* \) on an arbitrary basis element \( e_x \):
    \[
        (m m^*)(e_x) = m\big(m^*(e_x)\big) = m(e_x \otimes e_x) = e_x e_x = e_x.
    \]

    Because the linear operators \( m m^* \) and \( \text{Id} \) agree on
    every element of the basis \( \{e_x\}_{x \in X} \), by linearity they
    must be equal on the entire space \( C(X) \). Therefore:
    \[
        m m^* = \text{Id}.
    \]
\end{example}

% TODO: Is this necessary?
\begin{example}
    In particular, \( (\mathbb{C}^{n}, \psi_{n}) \) with \( \psi_{n}(e_{i}) = 1
    \) is a quantum set. The functional naturally extends to:
    \[
        \psi_n(x) = \sum_{i=1}^n x_i.
    \]

    To show that \(\psi_n\) is positive, we evaluate it on the
    element \(x^* x\). Using the orthogonality of the basis, we have:
    \[
        x^* x = \left( \sum_{i=1}^n \overline{x_i} e_i \right) \left(
        \sum_{j=1}^n x_j e_j \right) = \sum_{i,j=1}^n \overline{x_i}
        x_j e_i e_j = \sum_{i=1}^n \overline{x_i} x_i e_i =
        \sum_{i=1}^n |x_i|^2 e_i.
    \]
    Applying the functional yields:
    \[
        \psi_n(x^* x) = \psi_n\left( \sum_{i=1}^n |x_i|^2 e_i \right)
        = \sum_{i=1}^n |x_i|^2 \psi_n(e_i) = \sum_{i=1}^n |x_i|^2.
    \]
    Since \(|x_i|^2 \geq 0\) for all \(i\), the sum of these
    non-negative real numbers is also non-negative, meaning
    \(\psi_n(x^* x) \geq 0\). Thus, the functional is positive.

    To prove faithfulness, suppose \(\psi_n(x^* x) = 0\) for some \(x
    \in \mathbb{C}^n\). This implies:
    \[
        \sum_{i=1}^n |x_i|^2 = 0.
    \]
    Because each term \(|x_i|^2\) is strictly non-negative, the sum
    can only be zero if every individual component satisfies
    \(|x_i|^2 = 0\) for all \(i \in \{1, \dots, n\}\). This directly
    implies \(x_i = 0\) for all \(i\). Therefore, \(x = \sum_{i=1}^n
    0 \cdot e_i = 0\), establishing that \(\psi_n\) is faithful.

    The adjoint map \(m^*\colon \mathbb{C}^n
    \rightarrow \mathbb{C}^n \otimes \mathbb{C}^n\) acts on a basis
    element \(e_i\) by:
    \[
        m^*(e_i) = e_i \otimes e_i.
    \]

    We wish to show that the composition \(mm^*\colon \mathbb{C}^n
    \rightarrow \mathbb{C}^n\) is the identity operator. We check:
    \[
        (m m^*)(e_i) = m\big(m^*(e_i)\big) = m(e_i \otimes e_i) = e_i.
    \]
    Because the linear operators \(m m^* \) and \(\text{Id}\) agree
    on every element of the basis \(\{e_i\}_{i=1}^n\), by linearity
    they must be equal on the entire vector space \(\mathbb{C}^n\). Therefore:
    \[
        m m^* = \text{Id}.
    \]
\end{example}

% TODO: Why is there nTr instead of Tr?
\begin{example}
    For any \( n \in \mathbb{N}_{+} \), \( (M_{n}, n
    \operatorname{Tr}) \) is a quantum set, with \( \operatorname{Tr}
    \) being the trace. We will later show that this structure can be
    obtained as the crossed product \( (\mathbb{C}^{n}, \psi_{n}) \) and
    \( \hat{\mathbb{Z}}_{n} \) with the standard action, i.e.
    \[
        (M_{n}, n \operatorname{Tr}) \cong (\mathbb{C}^{n} \rtimes
        \hat{\mathbb{Z}}_{n}, \psi_{\rtimes}).
    \]
\end{example}

% TODO: give a better name
\begin{lemma}[Comultiplication for matrix algebras]
    In the context of \( (M_{n}, n \operatorname{Tr}) \), the
    comultiplication
    is given by the formula
    \[
        m^{*}(e_{ij}) = \sum_{k=1}^{n} e_{ik} \otimes e_{kj}
    \]
    with \( e_{ij} \) being the matrix with a single \( 1 \) in
    the \( i \)-th
    row and \( j \)-th column, and extended linearly.
\end{lemma}

% TODO: check this proof carefully
\begin{proof}
    This is simply a computational excercise. We will use the following
    well-known formulas:
    \begin{enumerate}
        \item \( e_{ab}e_{cd} = \delta_{bc}e_{ad} \),
        \item \( \left\langle e_{ij}, e_{kl} \right\rangle =
            \operatorname{Tr}(e_{ji}e_{kl}) = \delta_{ik}e_{jl} \),
    \end{enumerate}
    with
    \[
        \delta_{ij} =
        \begin{cases}
            1 & \text{if } i = j,\\
            0 & \text{otherwise}.
        \end{cases}
    \]
    In particular \( \left\{ e_{ij} \right\}_{i,j=1,\dots,n} \) form
    a orthonormal basis.

    Take arbitrary basis vectors \( e_{ab} \), \( e_{cd} \). Then
    \[
        m(e_{ab} \otimes e_{cd}) = e_{ab} e_{cd} = \delta_{bc}e_{ad},
    \]
    hence
    \begin{align*}
        \left\langle m(e_{ab} \otimes e_{cd}), e_{ij} \right\rangle
        &= \delta_{bc} \left\langle e_{ad}, e_{ij} \right\rangle \\
        &= \delta_{bc}\delta_{ai}\delta_{dj}.
    \end{align*}
    Now compute the inner product with the proposed adjoint:
    \begin{align*}
        \left\langle e_{ab} \otimes e_{cd}, \sum_{k=1}^{n} e_{ik}
        \otimes e_{kj} \right\rangle
        &= \sum_{k=1}^{n} \left\langle
        e_{ab}, e_{ik} \right\rangle
        \left\langle e_{cd}, e_{kj} \right\rangle \\
        &= \sum_{k=1}^{n}
        \delta_{ai}\delta_{bk}\delta_{ck}\delta_{dj} \\
        &= \delta_{ai}\delta_{bc}\delta_{dj}.
    \end{align*}
    The two expressions agree. Since the matrix unix form an orthonormal
    basis, this linearly extends to all matrices.
\end{proof}

\subsection{Crossed product quantum sets}

\textcite{tobolski26} have introduced the notion of crossed product quantum
sets. Suppose we are given a finite-dimensional \(C^{*}\)-algebra
\( \tilde{B}
\) along with a dual group action \( \hat{a}\colon \hat{G} \rightarrow
\operatorname{Aut}(\tilde{B}) \) where \( G \) is finite and
abelian. One can
then consider the crossed product
\[
    B \coloneqq \tilde{B} \rtimes_{\hat{\alpha}} \hat{G}.
\]

If \( \tilde{B} \) is equipped with a faithful functional \( \tilde{\psi} \)
in such a way as to make \( (\tilde{B}, \tilde{\psi}) \) a quantum
set, one can define a new faithful functional on \( B \) to make it a
quantum set, as described below.

\begin{definition}[Crossed product quantum set,
    \textcite[Definition~3.1]{tobolski26}]
    Let \( (\tilde{B}, \tilde{\psi}) \) be a quantum set and let \(
        \hat{\alpha}\colon \hat{G} \rightarrow \operatorname{Aut}(\tilde{B},
    \tilde{\psi}) \) be an action of the dual of a finite abelian
    group \( G \). The \emph{crossed product quantum set} is defined as
    \[
        (\tilde{B}, \tilde{\psi}) \rtimes_{\hat{\alpha}} \hat{G} \coloneqq
        \left( \tilde{B} \rtimes_{\hat{\alpha}} \hat{G},
        \psi_{\rtimes} \right),
    \]
    where \( \tilde{B} \rtimes_{\hat{\alpha}} \hat{G} \) is the
    crossed product \(C^{*}\)-algebra and \( \psi_{\rtimes} \) is the
    functional given by
    \[
        \psi_{\rtimes} \left( \sum_{\chi \in \hat{G}}
        b_{\chi}u_{\chi} \right) = n \tilde{\psi}(b_{1})
    \]
    with \( n = | \hat{G} | \).
\end{definition}

The proof that this definition is sensible, i.e. the newly defined
object is indeed a quantum set, can be seen in \textcite{tobolski26},
Proposition 3.4.

% TODO: Prove this lemma
\begin{lemma}
    Let \( \hat{\alpha}\colon \hat{\mathbb{Z}}_{n} \rightarrow
    \operatorname{A} \) be
    the action of cyclically shifting elements. Then
    \[
        (M_{n}, n \operatorname{Tr}) \cong (\mathbb{C}^{n} \rtimes_{\alpha}
        \hat{\mathbb{Z}}_{n}, \psi_{\rtimes}).
    \]
\end{lemma}

\section{Quantum graphs}

A simple graph is classically defined as a pair \( (V, E) \) with \( V \)
serving as the \emph{set of vertices} and \( E \subseteq [V]^{2} \)
serving as the \emph{set of edges}. Such a graph could be interpreted
as a matrix \( A = [a_{ij}] \in \{0, 1\}^{n \times n} \) with \(
\{v_{i}, v_{j}\} \in E \Leftrightarrow a_{ij} = 1 \); such a
representation is often called the \emph{adjacency matrix}. The key
property of such matrices is that they are invariant with respect to
entrywise multiplication. In turn we can use this property to
generalise graphs on sets to graphs on quantum sets.

The following definitions are due to \textcite{tobolski26}.
% TODO: Tell where this notion was first introduced
\begin{definition}[Quantum adjacency matrix,
    \textcite[Definition~2.5]{tobolski26}]
    Let \( (B, \psi) \) be a quantum set. A \emph{quantum adjacency
    matrix} on this quantum set is a linear map \( A\colon B \rightarrow B
    \) which is Schur-idempotent and *-preserving, by which we mean
    that it satisfies
    \[
        m(A \otimes A)m^{*} = A,
        \qquad
        A(b^{*}) = A(b)^{*}
        \quad \text{for all } b \in B.
    \]
\end{definition}

\begin{definition}[Quantum graph, \textcite[Definition~2.5]{tobolski26}]
    A \emph{quantum graph} on \( (B, \psi) \) is given by a quantum
    adjacency matrix on this quantum set.
\end{definition}

For all intents and purposes we will identify both notions and
simply call an
appropriate \( A\colon B \rightarrow B \) both a quantum adjancenty matrix and a
quantum graph.

\begin{remark}
    There are other equivalent definitions of a quantum graph, namely those
    based on
    \begin{itemize}
        \item an \emph{edge projection}, i.e. a projection \(
            \tilde{A} \in C(X) \otimes C(X)^{\operatorname{op}} \)
            \textcite{musto18},
        \item
    \end{itemize}
\end{remark}

\begin{definition}[\textcite[Definition~4.6{tobolski26}]]
    Let \( A \) be a quantum graph on a quantum set \( (B, \psi) \).
    We call the quantum graph
    \begin{itemize}
        \item \emph{loopfree} if \( m(A \otimes
            \operatorname{Id})m^{*} = 0, \)
        \item \emph{undirected} if \( A = A^{*} \), where the adjoint
            is taken with respect to the inner product induced by
            \( \psi \).
    \end{itemize}
    We say that a voltage quantum graph \( \tilde{A} = \left\{
    \tilde{A}_{g} \right\}_{g \in G} \) has these properties if the
    same equations hold for \( \sum_{g \in G} \tilde{A}_{g} \),
    respectively.
\end{definition}

\section{Isomorphisms and Quantum Isomorphisms of Quantum Sets}

\begin{definition}[Quantum isomorphisms,
    \textcite[Definition~2.8]{tobolski26}]
    Let \( (B_{1}, \psi_{1}) \) and \( (B_{2}, \psi_{2}) \) be
    quantum sets and
    let \( A_{1}\colon B_{1} \rightarrow B_{1} \) and \( A_{2}\colon B_{2}
    \rightarrow B_{2} \)
    be quantum graphs on \( (B_{1}, \psi_{1}) \), \( (B_{2},
    \psi_{2}) \), respectively.
    \begin{itemize}
        \item The two quantum sets are \emph{isomorphic} if there is
            a \( * \)-isomorphism
            \( \phi\colon B_{1} \rightarrow B_{2} \) such that \( \psi_{2}
            \phi = \psi_{1} \).
            We denote by \( \operatorname{Aut}(B, \psi) \) the group
            of quantum set
            automorphisms of a quantum set \( (B, \psi) \).
        \item The two quantum graphs are \emph{isomorphic} if there
            is an isomorpshism
            \( \phi\colon B_{1} \rightarrow B_{2} \) of the underlying
            quantum sets such that
            \( \phi A_{1} = A_{2} \phi \). We denote by \(
            \operatorname{Aut}(A) \) the
            group of quantum graph automorphisms of a quantum graph \( A \).
        \item The two quantum sets are \emph{quantum isomorphic}, written
            \( (B_{1}, \psi_{1}) \cong (B_{2}, \psi_{2}) \), if there is
            a finite-dimensional Hilbert space \( \mathcal{H} \) and a
            unital \( * \)-homomorphism
            \[
                \rho\colon B_{1} \rightarrow B_{2} \otimes B(\mathcal{H}),
            \]
            such that the map
            \[
                p\colon L^{2}(B_{1}) \otimes \mathcal{H} \rightarrow
                L^{2}(B_{2}) \otimes \mathcal{H}, \; a \otimes \xi
                \mapsto \rho(a)(1 \otimes \xi)
            \]
            is unitary as an operator between Hilbert spaces.
        \item The two quantum graphs are \emph{quantum isomorphic},
            written \( A_{1} \cong_{q} A_{2} \), if there is a
            quantum isomorphism \( \rho\colon B_{1} \rightarrow B_{2}
            \otimes B(\mathcal{H}) \) between the underlying quantum
            sets such that
            \[
                \rho A_{1} = (A_{2} \otimes
                \operatorname{Id}_{B(\mathcal{H})}) \rho
            \]
            holds.
    \end{itemize}
\end{definition}

\chapter{Quantum graphs over \( 3 \times 3 \) matrices}

The classification on \( 2 \times 2 \) matrices was done by
\textcite{matsuda22}.

\section{Gromada's theorem}

In the following section we will state...

The following theorem is due to \textcite{gromada22}.

\begin{theorem}[\textcite{gromada22}, Theorem 3.7]
    Consider \( m \in \left\{ 0, 1, \dots, n^{2} \right\} \). Quantum
    graphs on \( M_{n} \) with \( m \) quantum edges are in a
    one-to-one correspondence with \( m \)-dimensional subspaces \( V
    \subseteq \mathfrak{gl}_{n} \). The adjacency matrix is given by
    \( A_{V} = \sum_{s=1}^{m} \xi_{s} \otimes \xi_{s}^{*} \), where
    \( (\xi_{1}, \xi_{2}, \dots, \xi_{m}) \) is some orthonormal basis
    of \( V \).

    The quantum graph has no loops if and only if \( V \subseteq
    \mathfrak{sl}_{n} \).
    It is undirected if and only if \( V = V^{*} \).
\end{theorem}

\begin{proof}
    See Gromada's paper for reference.
\end{proof}

\begin{lemma}[\textcite{gromada22}, Lemma 3.8]
    Any vector subspace \( V \subseteq \mathfrak{gl}_{n} \) with \( V
    = V^{*} \) has a self-adjoint basis \( \left\{ \xi_{s} \right\}
    \), \( \xi_{s} = \xi_{s}^{*} \).
\end{lemma}

\begin{proof}
    We construct the basis by induction. Suppose we already have a
    basis \( \left\{ \xi_{s} \right\} \) of a subspace \( V_{1}
    \subseteq V \) such that \( \xi_{s} = \xi_{s}^{*} \). This means
    that \( V_{1} = V_{1}^{*} \). Let \( V_{2} \) be complement to \(
    V_{1} \). Take any \( \xi \in V_{2} \) and consider \( (\xi +
    \xi^{*})/2 \) and \( i(\xi^{*} - \xi)/2 \). Both elements are
    self-adjoint and elements of \( V_{2} \). They might or might not
    be linearly independent, but definitely at least one of them is
    non-zero. So, include both (if they are linearly independent) or
    one of them (if they are not) to our basis \( \left\{ \xi_{s}
    \right\} \) and repeat.
\end{proof}

\begin{remark}
    The following lemma tells us that any undirected graph on \( M_{n} \) can
    be thought of as being made from symmetric quantum edges only.
\end{remark}

\begin{remark}
    A complex vector space \( V = V^{*} \subseteq \mathfrak{gl}_{n}
    \) uniquely determines a real vector space spanned by the basis
    \( \left\{ \xi_{s} \right\} \) of self-adjoint elements \( \xi_{s} =
    \xi_{s}^{*} \). Indeed, it is easy to see that any other self-adjoint
    element \( \xi = \xi^{*} \) must be a real linear combination of
    \( \left\{ \xi_{s} \right\} \).
\end{remark}

\chapter{Conclusion}

blabla

\section{Further research}
one could try to classify graphs on \( M_n \) but it looks kinda tough ngl


%%%%% BIBLIOGRAFIA
\printbibliography

\end{document}
